\begin{proofbox}
	\textbf{\textcolor{CProof}{思路提示}}

	建立空间直角坐标系，先严格确定球心位置，再计算球心到任一顶点的距离，即得到外接球半径。

	\medskip
	\textbf{\textcolor{CProof}{正式推导}}
	\begin{description}[style=nextline, leftmargin=2.8em, labelsep=0.5em, itemsep=0.55em]
		\item[\textbf{步骤一（建立坐标系）：}]
		      取长方体一个顶点为原点 $O(0,0,0)$，过 $O$ 的三条棱分别沿 $x$ 轴、$y$ 轴、$z$ 轴正向，棱长分别为 $a,b,c$。

		      则长方体的顶点坐标可统一表示为
		      \[
			      (x,y,z),\quad x\in\{0,a\},\quad y\in\{0,b\},\quad z\in\{0,c\}.
		      \]
		      其中三个与 $O$ 相邻的顶点和一个对顶点可记为
		      \[
			      A(a,0,0),\quad B(0,b,0),\quad C(0,0,c),\quad D(a,b,c).
		      \]
		      \par\textbf{理由：}
		      这样建系后，长方体的长、宽、高分别对应三个坐标方向，距离计算最简单。

		\item[\textbf{步骤二（严格确定球心）：}]
		      设长方体外接球球心为 $P(x,y,z)$。因为外接球经过长方体所有顶点，所以
		      \[
			      |PO|=|PA|=|PB|=|PC|.
		      \]
		      由 $|PO|^2=|PA|^2$，得
		      \[
			      x^2+y^2+z^2=(x-a)^2+y^2+z^2,
		      \]
		      所以
		      \[
			      x=\frac{a}{2}.
		      \]
		      同理，由 $|PO|^2=|PB|^2$ 和 $|PO|^2=|PC|^2$，得
		      \[
			      y=\frac{b}{2},\qquad z=\frac{c}{2}.
		      \]
		      因此外接球球心为
		      \[
			      M\left(\frac{a}{2},\frac{b}{2},\frac{c}{2}\right).
		      \]
		      \par\textbf{理由：}
		      球心到球面上各顶点的距离相等。用等距关系可以直接推出球心坐标，避免只凭直观判断。

		\item[\textbf{步骤三（验证所有顶点都在同一球面上）：}]
		      任取长方体的一个顶点
		      \[
			      V(\varepsilon_1 a,\varepsilon_2 b,\varepsilon_3 c),
			      \qquad \varepsilon_1,\varepsilon_2,\varepsilon_3\in\{0,1\}.
		      \]
		      则
		      \[
			      |MV|^2
			      =
			      \left(\varepsilon_1a-\frac{a}{2}\right)^2
			      +
			      \left(\varepsilon_2b-\frac{b}{2}\right)^2
			      +
			      \left(\varepsilon_3c-\frac{c}{2}\right)^2.
		      \]
		      由于 $\varepsilon_i$ 只能取 $0$ 或 $1$，所以每一项的平方分别为
		      \[
			      \frac{a^2}{4},\quad \frac{b^2}{4},\quad \frac{c^2}{4}.
		      \]
		      因此
		      \[
			      |MV|^2=\frac{a^2+b^2+c^2}{4}.
		      \]
		      这说明 $M$ 到长方体所有顶点的距离都相等。

		      \par\textbf{理由：}
		      这一步说明以 $M$ 为球心、$|MV|$ 为半径的球确实经过长方体全部顶点。

		\item[\textbf{步骤四（计算半径）：}]
		      外接球半径为
		      \[
			      R=|MO|
			      =
			      \sqrt{
				      \left(\frac{a}{2}\right)^2
				      +
				      \left(\frac{b}{2}\right)^2
				      +
				      \left(\frac{c}{2}\right)^2
			      }
			      =
			      \frac{1}{2}\sqrt{a^2+b^2+c^2}.
		      \]
		      \par\textbf{理由：}
		      由空间两点间距离公式计算球心到顶点的距离。
	\end{description}

	\medskip
	\textbf{\textcolor{CProof}{结论回扣}}

	\textbf{因此，长方体外接球的球心在长方体中心，半径为}
	\[
		\boxed{R=\frac{1}{2}\sqrt{a^2+b^2+c^2}}.
	\]
\end{proofbox}